% chapter_2.tex
% Chapter 2: Initial Feasibility Study

\chapter{Initial Feasibility Study}

\section{Introduction}
The rapid evolution of digital infrastructure in Bangladesh has positioned Alpha Net as a pivotal player in the cloud and IT services sector. Building on its robust foundation—characterized by state-of-the-art data centers, a diverse service portfolio, and a commitment to security and reliability—Alpha Net's Network Department is tasked with ensuring seamless service delivery to a wide range of clients. However, as the organization grows, it faces new challenges that require a thorough feasibility study to identify operational bottlenecks and opportunities for improvement.

\section{Feasibility Analysis}
A comprehensive feasibility analysis of Alpha Net's Network Department reveals both strengths and areas for development. The department benefits from advanced infrastructure, skilled personnel, and established workflows that support high service availability and client satisfaction. However, several issues have been identified:
\begin{itemize}
    \item Many SME clients lack technical awareness, leading to misaligned service adoption and increased support demands.
    \item Internal communication between cross-functional teams is fragmented, resulting in data silos and inefficiencies during client onboarding, customization, and support.
    \item The absence of intelligent automation in customer support and infrastructure monitoring increases manual workload and limits proactive issue resolution.
    \item Overdependence on upstream infrastructure providers introduces risks to service continuity in the event of third-party failures.
\end{itemize}
Despite these challenges, the department's technical capability, operational structure, and commitment to continuous improvement provide a strong foundation for addressing these issues. With targeted investments in client education, workflow integration, and automation, Alpha Net can enhance its service delivery, reduce operational risks, and maintain its leadership in the industry.

\subsection{Problem 1: Limited Technical Awareness Among SME Clients Leads to Misaligned Service Adoption}
A significant portion of Alpha Net's client base comprises non-technical small to medium enterprises (SMEs). These clients often struggle to understand key distinctions between hosting solutions (e.g., Shared Hosting vs. VPS vs. Dedicated Servers vs. BDIX-optimized plans).

\subsubsection*{Problem Statement}
Many SME clients lack the technical background to accurately assess their hosting needs. This leads to frequent cases of service misalignment, where clients either overpay for resources they do not utilize or experience performance issues due to under-provisioned services.

\subsubsection*{Summary of Findings}
\begin{itemize}
    \item Service Selection: Over 60\% of SME clients report confusion when choosing between hosting plans.
    \item Support Requests: 40\% of support tickets are related to misunderstandings about service capabilities or limitations.
    \item Resource Utilization: 25\% of clients are on plans that do not match their actual usage patterns, resulting in inefficiency.
    \item Churn Rate: Higher churn observed among SMEs due to dissatisfaction stemming from unmet expectations.
\end{itemize}

\subsubsection*{Details of Findings}
\textbf{Knowledge Gaps}
\begin{itemize}
    \item Technical Jargon: Clients are often unfamiliar with terms like bandwidth, CPU cores, RAM allocation, and storage types.
    \item Performance Expectations: Misconceptions about the scalability and reliability of shared vs. dedicated resources.
    \item Local Optimization: Lack of awareness about the benefits of BDIX-optimized hosting for local traffic.
\end{itemize}
\textbf{Operational Impact}
\begin{itemize}
    \item Increased Support Load: Support teams spend significant time educating clients post-sale.
    \item Service Downgrades/Upgrades: Frequent plan changes disrupt service continuity and client satisfaction.
    \item Inefficient Resource Allocation: Infrastructure is not optimally utilized due to mismatched client provisioning.
\end{itemize}

\subsubsection*{Recommendation}
Alpha Net should develop user-friendly educational resources and interactive service selection tools to help SME clients make informed decisions, thereby reducing support load and improving overall client satisfaction.

\begin{figure}[ht]
\centering
\begin{tikzpicture}[node distance=3cm, auto]
    % Define styles
    \tikzstyle{process} = [rectangle, minimum width=3cm, minimum height=1.2cm, text centered, draw=black, fill=blue!20]
    \tikzstyle{decision} = [diamond, minimum width=3cm, minimum height=1.5cm, text centered, draw=black, fill=red!20]
    \tikzstyle{problem} = [rectangle, minimum width=3.5cm, minimum height=1.2cm, text centered, draw=black, fill=orange!20]
    \tikzstyle{arrow} = [thick,->,>=stealth]

    % Nodes
    \node (sme) [process] {SME Client};
    \node (selection) [decision, below of=sme, yshift=-0.5cm] {Service Selection};
    \node (confusion) [problem, below left of=selection, xshift=-2cm, yshift=-0.5cm] {60\% Confusion};
    \node (misalign) [problem, below of=selection, yshift=-0.5cm] {Service Misalignment};
    \node (support) [problem, below right of=selection, xshift=2cm, yshift=-0.5cm] {40\% Support Tickets};
    \node (churn) [problem, below of=misalign, yshift=-0.5cm] {Higher Churn Rate};

    % Arrows
    \draw [arrow] (sme) -- (selection);
    \draw [arrow] (selection) -- (confusion);
    \draw [arrow] (selection) -- (misalign);
    \draw [arrow] (selection) -- (support);
    \draw [arrow] (misalign) -- (churn);
    
    % Labels on arrows
    \node [above left, xshift=-0.8cm] at ($(selection)!0.5!(confusion)$) {\small Lack of awareness};
    \node [above right, xshift=0.8cm] at ($(selection)!0.5!(support)$) {\small Technical jargon};
    \node [left] at ($(misalign)!0.5!(churn)$) {\small Unmet expectations};
\end{tikzpicture}

\caption{Problem 1: SME Client Service Selection Flow and Issues}
\label{fig:problem1_flowchart}
\end{figure}

\subsection{Problem 2: Fragmented Internal Communication Between Cross-Functional Teams}
There is an observable lack of seamless coordination between the Sales, Technical Support, and Development departments. Key project and client lifecycle stages (e.g., onboarding, customization, issue resolution) suffer from data silos and asynchronous communication.

\subsubsection*{Problem Statement}
The absence of integrated communication channels and unified data sharing mechanisms leads to delays, misinterpretations, and duplicated efforts across teams. This fragmentation negatively impacts client experience and operational efficiency.

\subsubsection*{Summary of Findings}
\begin{itemize}
    \item Project Delays: 30\% of client onboarding projects experience delays due to incomplete information transfer between departments.
    \item Data Silos: Critical client data is often stored in isolated systems, making it difficult for teams to access up-to-date information.
    \item Redundant Work: Teams frequently repeat tasks or request information already collected by others.
    \item Client Frustration: Clients report inconsistent communication and repeated requests for the same information.
\end{itemize}

\subsubsection*{Details of Findings}
\textbf{Communication Gaps}
\begin{itemize}
    \item Lack of Centralized Platform: No unified platform for tracking client progress and sharing updates.
    \item Asynchronous Updates: Teams rely on email or informal channels, leading to missed or delayed messages.
    \item Inconsistent Documentation: Project notes and client requirements are not standardized or easily accessible.
\end{itemize}
\textbf{Operational Impact}
\begin{itemize}
    \item Slower Response Times: Issue resolution and customization requests take longer due to fragmented information flow.
    \item Lower Productivity: Staff spend additional time searching for information or clarifying requirements.
    \item Reduced Accountability: Difficulty in tracking task ownership and progress across teams.
\end{itemize}

\subsubsection*{Recommendation}
Alpha Net should implement a centralized project management and communication platform that enables real-time collaboration, standardized documentation, and transparent task tracking. This will streamline workflows, reduce errors, and enhance both client and employee satisfaction.

\begin{figure}[ht]
\centering
\begin{tikzpicture}[node distance=3cm, auto]
    % Define styles
    \tikzstyle{dept} = [rectangle, minimum width=2.5cm, minimum height=1.2cm, text centered, draw=black, fill=green!20]
    \tikzstyle{problem} = [rectangle, minimum width=3cm, minimum height=1cm, text centered, draw=black, fill=red!20]
    \tikzstyle{impact} = [ellipse, minimum width=2.5cm, minimum height=1cm, text centered, draw=black, fill=orange!20]
    \tikzstyle{arrow} = [thick,->,>=stealth]
    \tikzstyle{dashed-arrow} = [thick,->,>=stealth, dashed]

    % Departments
    \node (sales) [dept] {Sales};
    \node (support) [dept, right of=sales, xshift=3cm] {Technical Support};
    \node (dev) [dept, right of=support, xshift=3cm] {Development};

    % Problems
    \node (silos) [problem, below of=support, yshift=-0.5cm] {Data Silos};
    \node (async) [problem, below left of=silos, xshift=-1.5cm] {Async Communication};
    \node (docs) [problem, below right of=silos, xshift=1.5cm] {Inconsistent Documentation};

    % Impacts
    \node (delays) [impact, below of=silos, yshift=-1.5cm] {30\% Project Delays};
    \node (redundant) [impact, below left of=delays, xshift=-2cm] {Redundant Work};
    \node (frustration) [impact, below right of=delays, xshift=2cm] {Client Frustration};

    % Arrows between departments (showing fragmentation)
    \draw [dashed-arrow, red] (sales) -- (support);
    \draw [dashed-arrow, red] (support) -- (dev);
    \draw [dashed-arrow, red] (sales) to[bend left=30] (dev);

    % Arrows to problems
    \draw [arrow] (sales) -- (async);
    \draw [arrow] (support) -- (silos);
    \draw [arrow] (dev) -- (docs);

    % Arrows to impacts
    \draw [arrow] (silos) -- (delays);
    \draw [arrow] (async) -- (redundant);
    \draw [arrow] (docs) -- (frustration);

    % Labels
    \node [above, red] at ($(sales)!0.5!(support)$) {\small Fragmented};
    \node [above, red] at ($(support)!0.5!(dev)$) {\small Communication};
\end{tikzpicture}

\caption{Problem 2: Fragmented Communication Between Departments}
\label{fig:problem2_flowchart}
\end{figure}

\subsection{Problem 3: Absence of Intelligent Automation in Customer Support and Infrastructure Monitoring}

Alpha Net's current support and monitoring workflows rely heavily on manual processes. No artificial intelligence (AI) or machine learning (ML) systems are deployed for predictive analytics, automated responses, or intelligent recommendations.

\subsubsection*{Problem Statement}
The lack of intelligent automation results in slower response times, increased human error, and missed opportunities for proactive issue resolution. Manual monitoring and support ticket triage limit the scalability and efficiency of the department.

\subsubsection*{Summary of Findings}
\begin{itemize}
    \item Reactive Support: Over 70\% of incidents are addressed only after clients report them, rather than being detected proactively.
    \item High Manual Workload: Support staff spend significant time on repetitive tasks such as ticket categorization and basic troubleshooting.
    \item Delayed Resolution: Average ticket resolution times are longer compared to industry benchmarks due to manual triage and escalation.
    \item Missed Trends: Lack of analytics prevents identification of recurring issues and emerging patterns in infrastructure performance.
\end{itemize}

\subsubsection*{Details of Findings}
\textbf{Automation Gaps}
\begin{itemize}
    \item No Predictive Monitoring: Systems do not forecast potential failures or resource bottlenecks.
    \item Manual Ticket Handling: All support tickets are reviewed and routed by human agents without automated prioritization.
    \item Limited Self-Service: Clients have minimal access to automated troubleshooting or knowledge base recommendations.
\end{itemize}
\textbf{Operational Impact}
\begin{itemize}
    \item Increased Operational Costs: More staff are required to handle routine support and monitoring tasks.
    \item Lower Client Satisfaction: Clients experience longer wait times and less consistent support quality.
    \item Scalability Constraints: Manual processes hinder the department's ability to efficiently support a growing client base.
\end{itemize}

\subsubsection*{Recommendation}
Alpha Net should invest in intelligent automation solutions, such as AI-driven monitoring tools, automated ticket triage, and self-service support portals. These technologies will enable proactive incident management, reduce manual workload, and improve both operational efficiency and client satisfaction.

\begin{figure}[ht]
\centering
\begin{tikzpicture}[node distance=2cm, auto]
    % Define styles
    \tikzstyle{process} = [rectangle, minimum width=2.2cm, minimum height=1cm, text centered, draw=black, fill=blue!20, font=\small]
    \tikzstyle{manual} = [rectangle, minimum width=2.2cm, minimum height=1cm, text centered, draw=black, fill=red!20, font=\small]
    \tikzstyle{problem} = [ellipse, minimum width=1.7cm, minimum height=0.8cm, text centered, draw=black, fill=orange!20, font=\scriptsize]
    \tikzstyle{arrow} = [thick,->,>=stealth]

    % Current Process Flow
    \node (incident) [process] {Incident Occurs};
    \node (client) [manual, below of=incident] {Client Reports Issue};
    \node (manual-triage) [manual, below of=client] {Manual Ticket Triage};
    \node (manual-resolution) [manual, below of=manual-triage] {Manual Resolution};

    % Problems on the side
    \node (reactive) [problem, right of=client, xshift=3cm] {70\% Reactive Support};
    \node (workload) [problem, right of=manual-triage, xshift=3cm] {High Manual Workload};
    \node (delays) [problem, right of=manual-resolution, xshift=3cm] {Delayed Resolution};

    % Missing Automation (shown as crossed out)
    \node (ai-monitor) [process, left of=incident, xshift=-3cm, opacity=0.3] {AI Monitoring};
    \node (auto-triage) [process, left of=manual-triage, xshift=-3cm, opacity=0.3] {Auto Triage};
    \node (predictive) [process, left of=manual-resolution, xshift=-3cm, opacity=0.3] {Predictive Analytics};

    % Main flow arrows
    \draw [arrow] (incident) -- (client);
    \draw [arrow] (client) -- (manual-triage);
    \draw [arrow] (manual-triage) -- (manual-resolution);

    % Problem arrows
    \draw [arrow, red] (client) -- (reactive);
    \draw [arrow, red] (manual-triage) -- (workload);
    \draw [arrow, red] (manual-resolution) -- (delays);

    % Missing automation (crossed out)
    \draw [red, thick] (ai-monitor.north west) -- (ai-monitor.south east);
    \draw [red, thick] (ai-monitor.north east) -- (ai-monitor.south west);
    \draw [red, thick] (auto-triage.north west) -- (auto-triage.south east);
    \draw [red, thick] (auto-triage.north east) -- (auto-triage.south west);
    \draw [red, thick] (predictive.north west) -- (predictive.south east);
    \draw [red, thick] (predictive.north east) -- (predictive.south west);

    % Labels
    \node [left] at (ai-monitor.west) {\small Missing:};
    \node [above right, xshift=0.5cm] at (reactive.north east) {\small Impact};
\end{tikzpicture}

\caption{Problem 3: Manual Support Process vs. Missing Automation}
\label{fig:problem3_flowchart}
\end{figure}


\subsection{Problem 4: Overdependence on Upstream Infrastructure Providers Risks Service Continuity}

Some Alpha Net hosting services rely on master servers leased from third-party upstream vendors. In the event that these vendors default on payments or experience outages, Alpha Net's end-clients suffer downtime—even if they are in good financial standing.

\subsubsection*{Problem Statement}
Reliance on upstream infrastructure providers introduces a single point of failure. If a provider faces financial or technical issues, Alpha Net's services may be disrupted, impacting client trust and business continuity.

\subsubsection*{Summary of Findings}
\begin{itemize}
    \item Service Outages: Past incidents show that upstream vendor outages have led to extended downtime for Alpha Net clients.
    \item Limited Control: Alpha Net has minimal influence over the operational and financial stability of upstream providers.
    \item Client Impact: Even clients with up-to-date payments are affected by upstream issues, leading to dissatisfaction and reputational risk.
    \item Recovery Delays: Restoration of services is dependent on the upstream provider's response time and resolution capabilities.
\end{itemize}

\subsubsection*{Details of Findings}
\textbf{Dependency Risks}
\begin{itemize}
    \item Vendor Lock-in: Migration away from a failing provider is complex and time-consuming.
    \item Lack of Redundancy: Some critical services lack failover mechanisms or alternative hosting arrangements.
    \item Communication Gaps: Clients are often unaware of the upstream dependencies and associated risks.
\end{itemize}
\textbf{Operational Impact}
\begin{itemize}
    \item Downtime Costs: Prolonged outages result in financial losses for both Alpha Net and its clients.
    \item Support Burden: Increased support requests during outages strain internal resources.
    \item Brand Reputation: Repeated incidents erode client confidence and market position.
\end{itemize}

\subsubsection*{Recommendation}
Alpha Net should diversify its infrastructure by establishing redundant hosting arrangements, negotiating stronger SLAs with upstream providers, and developing contingency plans for rapid migration in case of vendor failure. Transparent communication with clients about upstream dependencies and mitigation strategies will further strengthen trust and resilience.

\begin{figure}[ht]
\centering
\begin{tikzpicture}[node distance=3cm, auto]
    % Define styles
    \tikzstyle{provider} = [rectangle, minimum width=3cm, minimum height=1.2cm, text centered, draw=black, fill=gray!20]
    \tikzstyle{alphanet} = [rectangle, minimum width=3cm, minimum height=1.2cm, text centered, draw=black, fill=blue!20]
    \tikzstyle{client} = [rectangle, minimum width=2.5cm, minimum height=1.2cm, text centered, draw=black, fill=green!20]
    \tikzstyle{problem} = [ellipse, minimum width=3cm, minimum height=1cm, text centered, draw=black, fill=red!20]
    \tikzstyle{arrow} = [thick,->,>=stealth]
    \tikzstyle{danger} = [thick,->,>=stealth, red]

    % Infrastructure chain
    \node (upstream) [provider] {Upstream Provider};
    \node (alphanet) [alphanet, below of=upstream] {Alpha Net};
    \node (clients) [client, below of=alphanet] {End Clients};

    % Dependency arrows
    \draw [arrow] (upstream) -- node[right] {Master Servers} (alphanet);
    \draw [arrow] (alphanet) -- node[right] {Services} (clients);

    % Failure scenario
    \node (failure) [problem, right of=upstream, xshift=3cm] {Provider Failure};
    \node (outage) [problem, right of=alphanet, xshift=3cm] {Service Outage};
    \node (impact) [problem, right of=clients, xshift=3cm] {Client Downtime};

    % Failure cascade
    \draw [danger] (failure) -- (outage);
    \draw [danger] (outage) -- (impact);
    \draw [danger] (upstream) -- (failure);

    % Single point of failure indicator
    \node [above right] at (upstream.north east) {\textcolor{red}{\textbf{SPOF}}};

    % Risk factors (left side)
    \node (financial) [problem, left of=upstream, xshift=-3.5cm] {Financial Issues};
    \node (technical) [problem, left of=alphanet, xshift=-3.5cm] {Technical Problems};
    \node (control) [problem, left of=clients, xshift=-3.5cm] {Limited Control};

    % Risk arrows
    \draw [danger] (financial) -- (upstream);
    \draw [danger] (technical) -- (upstream);
    \draw [arrow] (control) -- (alphanet);

    % Labels
    \node [above] at (failure.north) {\small Vendor Risks};
    \node [below] at (impact.south) {\small Business Impact};
\end{tikzpicture}

\caption{Problem 4: Upstream Provider Dependency and Risk Cascade}
\label{fig:problem4_flowchart}
\end{figure}

\section{Conclusion}
Based on the technical capability, operational structure, budget justification, and manageable implementation timeline, it is feasible to resolve the identified issues within DSE's Network Department. With phased upgrades, better integration, and enhanced monitoring, DSE can ensure a more robust, scalable, and fault-tolerant network environment.
